\pagebreak

{
\thispagestyle{empty}

\mbox{}

\vspace{7.5cm}

\footnotesize
\noindent{}\emph{\noindent{}Em meio aos ofícios da Inquisição da Igreja Católica em 1492, em Portugal e Espanha, contam as testemunhas que uma pessoa judia, se negando a converter-se à força ao catolicismo, é levada à fogueira. Junto a ela são colocados livros, folhas sagradas, pergaminhos e rolos da Torá para que sejam consumidos pelas chamas, junto à sua carne. A história nos conta que, no momento em que as primeiras labaredas incendiaram os papéis, as finas letras timbradas do alfabeto hebraico desprenderam-se, subindo em direção aos céus.\footnote{Segundo as histórias de Rabi Uri de Strelisk: ``As miríades de letras da Torá correspondem às miríades de almas judaicas. Se faltar uma letra no rolo da Torá, a Torá não será válida; se faltar uma alma na união judaica, a \emph{Shechiná} não há de pairar sobre ela. Como as letras, também as almas devem unir-se e formar uma só união. Por que, porém, se proíbe que uma letra da Torá toque a outra? Porque toda alma judaica deve ter sua hora em que possa ficar a sós em conexão divina''. Em: \textsc{buber}, \emph{História do Rabi}, p.\,464.} Deste registro flamejante de uma destruição que dispersa tudo --- coisas, letras, corpos, histórias, almas, espaços, lugares, tempos ---, desses múltiplos estilhaços --- de livros e de carnes --- outra coisa pode (re)nascer. Que ``deste reduzir às cinzas nasça o seu próprio dar à luz''.\footnote{\textsc{didi-huberman}, \emph{Esparsas: Viagens aos papéis do gueto de Varsóvia}, p.\,126.}}
}

% Comentário da editora: Essa nota não coube bem aqui… vi que ela também já aparece mais pra frente, no segundo capítulo.
% \footnote{\emph{Shechiná} é uma palavra hebraica que se refere, de forma resumida, à figura de Deus no feminino e sua habitação e presença no plano terrestre.}

\pagebreak

\chapter*{Epílogo\smallskip\subtitulo{Fagulhas}}
\addcontentsline{toc}{chapter}{Epílogo}

% Comentário da editora: Está estranho o pulo da Inquisição do gueto de Vilnius. Acho que ficou muito radical. Vi que depois volta-se a falar sobre ambos, fazendo a ligação, mas o texto não está sustentando uma tese. Tentei jogar pra didascália, pra ver se melhora. De toda forma, eu reescreveria esse capítulo com atenção.

% \noindent{}Em meio aos ofícios da Inquisição da Igreja Católica em 1492, em
% Portugal e Espanha, contam as testemunhas que uma pessoa judia, se
% negando a converter-se à força ao catolicismo, é levada à fogueira.
% Junto a ela são colocados livros, folhas sagradas, pergaminhos e rolos
% da Torá para que sejam consumidos pelas chamas, junto à sua
% carne. A história nos conta que, no momento em que as primeiras
% labaredas incendiaram os papéis, as finas letras timbradas do alfabeto
% hebraico desprenderam-se, subindo em direção aos céus.\footnote{Segundo
%   as histórias de Rabi Uri de Strelisk: ``As miríades de letras da
%   Torá correspondem às miríades de almas judaicas. Se faltar uma letra
%   no rolo da Torá, a Torá não será válida; se faltar uma alma na união
%   judaica, a Shechiná¹ não há de pairar sobre ela. Como as letras,
%   também as almas devem unir-se e formar uma só união. Por que, porém,
%   se proíbe que uma letra da Torá toque a outra? Porque toda alma
%   judaica deve ter sua hora em que possa ficar a sós em conexão
%   divina''. \emph{Shechiná} é uma palavra hebraica que se refere, de forma
%   resumida, à figura de Deus no feminino e sua habitação e presença no
%   plano terrestre. \textsc{buber}, \emph{História do Rabi}, p.\,464.} Deste registro
% flamejante de uma destruição que dispersa tudo --- coisas, letras,
% corpos, histórias, almas, espaços, lugares, tempos ---, desses múltiplos
% estilhaços --- de livros e de carnes --- outra coisa pode (re)nascer.
% Que ``deste reduzir às cinzas nasça o seu próprio dar à
% luz''\footnote{\textsc{didi-huberman}, \emph{Esparsas: Viagens aos papéis do
%   gueto de Varsóvia}, p.\,126.}.

\setlength{\epigraphwidth}{.65\textwidth}
\begin{epigraphs} 
\qitem{E vos digo: é preciso ter ainda caos dentro de si, para poder dar à luz uma estrela dançarina. E vos digo: ainda há caos dentro de vós.}
{
\textit{\textsc{taz} --- Zona Autônoma Temporária}\\%\smallskip
\textsc{hakim bey}
} 
\end{epigraphs}

% \begin{quote}
% E vos digo: é preciso ter ainda caos dentro de si, para poder dar à luz
% uma estrela dançarina. E vos digo: ainda há caos dentro de
% vós\footnote{\textsc{bey}, \emph{\textsc{taz} --- Zona Autônoma Temporária}, p.\,7.}.
% \end{quote}


% \begin{center}\adforn{49}\end{center}

% \noindent{}Disponho três folhas sulfite, uma ao lado da outra, sobre a mesa de
% jantar da sala de minha casa, rentes à beirada da tábua de madeira, em
% uma falha tentativa de criar um fundo branco. Sobre as folhas estão dois
% livros e um caderno. Registro esse primeiro passo meticulosamente
% realizado. Em seguida posiciono a tão esperada fotografia preta e
% branca, com alguns tons amarronzados, sobre os livros, o caderno, as
% folhas e a mesa. Procuro por algo que chame minha consciência desde uma
% certa distância --- temporal e espacial ---, uma ética de memória
% situada para além do binômio vida-morte\footnote{O binômio vida-morte se
%   estabelece no momento em que se lê as memórias trágicas do Holocausto
%   (e demais tragédias --- neste caso em específico, tragédias judaicas)
%   unicamente pela ótica da ``sobrevivência'' ou da ``não
%   sobrevivência''. Neste tipo de leitura, nós --- judeus pós-Segunda
%   Guerra Mundial (aqueles que sobreviveram --- nossa geração, nossos
%   pais e nossos filhos) --- nos colocamos, sem perceber ou até mesmo sem
%   querer, dentro de um espectro óptico traumático em que as memórias do
%   Holocausto cotidianamente nos atormentam e assombram. Afinal, não são
%   aqueles que sobreviveram que revisitam psiquicamente as situações dos
%   guetos, dos campos de concentração e de extermínio? Não somos nós,
%   sobreviventes, aqueles que pesquisam, debatem, choram e buscam ---
%   mesmo sem fim aparente --- respostas para explicar tais tragédias?
%   Neste texto procura-se entender que relembrar nossas dores e expor
%   nossas fragilidades é defender a memória, e assim, a vida.}.

\noindent{}Estou diante do retrato. Nele uma jovem mulher posa com sua cabeça
levemente inclinada, exibindo por completo a lateral esquerda de seu
afinado rosto. De bochechas rosadas e olhos esverdeados, seus cabelos de
tons castanhos escurecidos criam um contraste com sua pele clara
destacada contra o acinzentado do plano de fundo da imagem.

A mulher retratada é a minha trisavó Dvora Najzsteter\footnote{Mãe
  biológica de meu bisavô Aba (Alberto) Najzsteter \textsc{z''l}, que por
  sua vez é pai de minha avó materna Sarah Dora Najzsteter Haber
  \textsc{z''l} (1941--2021).}, nascida em algum vilarejo nas proximidades
da cidade de Lublin, na Polônia, e que agora, já em minhas mãos,
lança-se em uma superfície fotográfica fosca que, apesar do ligeiro
sorriso, exprime a crueldade de uma história de que pouco ou nada
saberei. Surge uma série de questionamentos. Não seria a fotografia a
aposta de um testemunho? Não seria a fotografia a possibilidade de
contato e de distância ao mesmo tempo? Cabe a nós não ter medo de olhar.

Neste lapso de tempo congelado, que data entre as décadas de 1930 e
1940, Dvora veste um blazer azul-petróleo com seis finas linhas brancas
nas golas, onde se destaca um broche prateado. Seus cabelos escuros,
curtos e ligeiramente ondulados, são fixados com esmero por finos
grampos de tonalidade preta, que poderiam passar despercebidos a um
olhar menos cuidadoso. Seus olhos turvos se desviam da lente fotográfica
em um gesto aparentemente proposital de desencontro. 

Será que de todas as coisas que poderiam estar olhando para ela agora,
ela imaginou por algum instante que seria eu?

No canto direito da fotografia, pelo olhar do espectador, está um
pequeno selo circular e timbrado que ocupa, aproximadamente, um quinto da
folha impressa. No centro se lê a palavra México, acompanhada dos
escritos: \textsc{foto herman e sta. maria de la redonda}, 136. Esse endereço não
existe mais. No verso do retrato, abrindo alas em meio aos vestígios do
tempo, encontra-se um pequeno texto, desejo de testemunho lido em
iídiche\footnote{O iídiche é uma língua judaico-germânica
  falada por judeus provenientes da Europa Central e Oriental (judeus
  ashkenazitas), uma mescla entre o alto-alemão do século \textsc{xiv} e
  elementos do hebraico e das línguas eslavas. Atualmente, apesar do
  grande número de falantes do iídiche serem judeus de vertentes
  ortodoxas e ultraortodoxas, o iídiche não é, em si, uma língua
  ortodoxa, mas sim uma língua diaspórica herdada e atravessada de
  espiritualidade e de tradição.} e escrito em letras hebraicas legíveis
e bem centralizadas na página, uma comprovação e, ao mesmo tempo, uma
palavra de adeus\footnote{Dvora fugiu da Polônia em torno do ano de
  1925, deixando o seu único filho --- órfão de pai --- aos cuidados de
  seus antigos sogros Dawid Najzsteter e Ruchla Lea Najzsteter. Ambos,
  Dawid e Ruchla Lea \textsc{z''l}, morreram em 1939 durante o início dos
  bombardeios nazistas nas cercanias da cidade de Lublin, na Polônia.}:
``Essa é uma fotografia minha de lembrança para os meus queridos
filhos. A tua mãe, Dvora''.

Sobre estes escritos --- ainda no verso do retrato --- encontra-se um
outro texto. Possivelmente uma primeira tentativa de aproximação entre a
autora, o papel e a caneta, mas pouco legível. As frases estão riscadas
e as letras se mesclam entre os traços azulados da caneta esferográfica.
As palavras parecem não possuir nem começo e nem fim.

% Comentário da editora: A preparadora trocou "este risco", que também não entendi a utilidade, por "os garranchos". São de fato garranchos? Não me incomoda, mas soa um pouco pejorativo pra um parágrafo que fala tanto sobre algo tão íntimo e pessoal. Também não entendi o trecho que segue, acho que precisa ser melhor explicado. A relação (sentimento, gancho com a pesquisa) com a fotografia da avó não está clara.

Os garranchos contrastam com a tonalidade já amarelada do suporte
fotográfico fosco, que me faz refletir. Penso sobre o apagamento da
lembrança, feito pela mão da própria autora --- e que ainda nestas
condições são lidas, desobedecendo o destino de uma morte prematura. São
gestos que não terminam, porque persistem no tempo. Como diz o poema de
Sutzkever: ``Talvez essas palavras também/\,Vão durar {[}\ldots{}{]}/\,Talvez essas palavras vão permanecer''\footnote{\textsc{sutzkever}
  \emph{apud} \textsc{didi-huberman}, \emph{op.\,cit.}, p.\,128.}. Já a segunda questão, o risco
da escrita por si só --- do ato de escrever letras que comprovam a
existência de um corpo, de uma pessoa em uma página. Uma mão, atrelada a
um corpo, escreveu essas palavras, que juntas se tornam frases,
finalmente são lidas enquanto desejo. \emph{Seria este o ritmo do fim?}
Entre essas promessas ardentes, prontas para serem consumidas pelas
chamas, estão também os papéis da Inquisição.

Ao olhar novamente percebo que, de minha trisavó Dvora, me restam
vestígios de sua caligrafia --- e assim uma terceira questão sobre o
risco surge\footnote{Proposta de risco como consequência.}:

% Comentário da editora: Eu acho que o iídiche deveria vir na nota, já que sempre priorizamos a acessibilidade do público (em português).

\begin{quote}
\emph{Ot di tokhike sakone tsu bashteyn oyf shraybn verter oyf yidish, vi di
vos shteyen ayngekritst in matn papir, vos ikh hob do far di oygn. Di-o
verter aleyn zaynen der konkreter bavayz fun iberlebn di natsishe
khayoles, fun firn a yidish lebn in der umgezetslekhkayt --- azoy az zey
zoln blaybn oyf yener zayt, bay der grenets tsvishn lebn un gezets. Azoy
arum az, in yenem moment, mer vi a shverd, antplekt zikh dos yidish fun
mayn elter-elter-bobe vi a pantser fun vidershtand kegn a
lingvistik-rasisher zoyberung, vos iz adurkhgefirt gevorn in eyrope. Dos
vayl, vi di yidish-meksikanishe shrayberin Myriam Moscona undz dermont:
\emph{Azoy vi mentshn, zaynen di shprakhn oykh farshikt gevorn in di
gazkamern}}\footnote{Este risco inerente de continuar a escrever
  palavras em iídiche como as que estão gravadas na folha fosca
  que agora tenho diante de meus olhos. Essas palavras são, por si só, a
  prova concreta de sobrevivência às tropas nazistas, à vida judaica na
  ilegalidade --- fazendo com que ela se mantenha no além, no limiar da
  vida e da Lei. E assim, naquele momento, mais do que uma espada, o
  iídiche de minha trisavó se mostra enquanto um escudo de
  resistência a uma higienização linguístico-racial em curso na Europa.
  Isso porque, como nos lembra a escritora judia-mexicana Myriam
  Moscona, ``assim como as pessoas, as línguas também foram
  enviadas às câmaras de gás''. \textsc{moscona}, \emph{Palabra del autor}.}
\end{quote}

Escrever em iídiche, em um gesto desesperado que vai de encontro
à fotografia, afugenta e retorna a mim. Como uma
resposta-eco\footnote{Trabalhar a experiência do eco como uma
  ressonância¹, algo que vai de encontro e que retorna em um duplo
  movimento de apropriação e distância. Pensar a ressonância enquanto tempo-limiar de retorno das
  perguntas que são colocadas no mundo. Ou seja, entre a primeira e a
  segunda aparição do signo, ou melhor, entre a ida e a volta de uma
  pergunta e/\,ou palavra ecoada, corre o tempo.}. Dessas memórias
permanece a \emph{nahalá}\footnote{Ritual judaico de reza e lembrança
  anual de entes queridos falecidos.}. E destas histórias de destruições
entoamos o \emph{Kadish}\footnote{Reza judaica de lembrança e elevação
  de alma e de espírito de entes falecidos em que há a santificação e a
  glorificação do nome de Deus.}. ``A nossa história não será em
vão''\footnote{Registro do Arquivo Ringelblum, acervo pertencente ao
  grupo judaico clandestino \emph{Oyneg Shabes} (em português, ``a
  alegria do \emph{shabat}''), que atuava coletando e registrando a vida
  comum e comunitária dentro do gueto de Varsóvia, entre os anos de
  1939--1943.}, suplicam todos aqueles judeus e judias que sobreviveram
ao desastre da vida judaica europeia entre as décadas de 1930 e 1940, e
continuam: ``todos saberão o que foi feito com os judeus na metade do século \textsc{xx}''\footnote{Registro do Arquivo
  Ringelblum.}. O iídiche, enquanto escudo, regressa para colocar
em xeque --- e assim, multiplicar --- tudo aquilo que por muito tempo
virou cinzas\footnote{A sobrevivência de uma língua diaspórica (como o
  iídiche e/\,ou o ladino) surge enquanto resposta às
  chamas flamejantes das fogueiras inquisitoriais e dos fornos
  crematórios dos campos de extermínio. São as próprias fagulhas das
  cinzas, esparsas partículas de papéis e de corpos que, lançadas para o
  ar e espalhadas pelo vento, se estabelecem e consolidam
  territorialidades próprias, dando sentidos de continuidade. Tanto o iídiche quanto o ladino são línguas
  judaicas diaspóricas. A primeira se constrói por meio de mesclas entre
  o alemão, o hebraico e alguns elementos de línguas eslavas. E a
  segunda por meio de mesclas entre o português e o espanhol medieval, o
  italiano e o hebraico. O iídiche é o idioma comum falado por
  judeus ashkenazitas (judeus provenientes da Europa Central e do
  Leste Europeu) enquanto o ladino é o idioma falado por judeus
  sefaraditas (judeus provenientes de Portugal e Espanha).}.

  % Comentário da editora: É só "metade" mesmo? Tem "primeira" ou algo assim? Seria necessário colocar [primeira] assim, entre colchetes, para indicar algo? Também acho que até dá pra entender o que "coloca em xeque" (que seria a própria validade da língua e da cultura iídiche), mas talvez faça mais sentido excluir essa frase. Isso foi algo que a preparadora apontou. Ou então reescrever essa parte.

Toda sobrevivência --- seja ela construída pelos seus maiores ou pelos
seus menores gestos --- é uma espécie de contestação histórica. Toda
sobrevivência é um desejo pulsante e latente de escrever\footnote{Para a
  maioria das pessoas judias perseguidas (mortas em bombardeios, dentro
  dos guetos, em campos de extermínio ou em valas comuns) não houve
  milagre. Suas vidas não foram salvas, para alguns apenas as suas
  palavras foram e, assim, persistem no tempo fazendo com que suas vidas
  e histórias sejam lidas ainda hoje, depois de suas mortes.}, de
registrar essa breve passagem, de fazer parar a roda da História. Assim
como os levantes, como diria o escritor e filósofo Hakim Bey\footnote{Em
  seu livro \emph{\textsc{taz} --- Zonas Temporárias Autônomas}, o escritor e
  filósofo anarquista Hakim Bey descreve a propagação de espaços
  autônomos temporários enquanto táticas de resistência e de
  esvaziamento de poder. Não seria o próprio ato de resistência do grupo
  clandestino judaico \emph{Oyneg Shabes} um espaço autônomo? Não seriam
  as páginas escritas e guardadas, para posteriormente serem resgatadas,
  criadoras de espacialidades temporárias? São possibilidades de outros
  mundos, e assim persistem no tempo.}, as tentativas de sobrevida são
insurreições que surgem para além do Tempo e violam a lei da História.
Os levantes --- em seus maiores ou em seus menores gestos --- e a
libertação de uma territorialidade --- em qualquer escala ---
possibilitam movimentos para fora e para além da \emph{espiral
hegeliana} de progresso. É dessa maneira que entendo as ligeiras
palavras em iídiche de minha trisavó --- um ato de registrar sua
breve passagem e gravar seus anseios de adeus e de amor.

Assim como nos contos talmúdicos, em que Deus criou o mundo (e
todas as coisas que ele contém) por meio da codificação das 22 letras
dispostas do alfabeto hebraico, o iídiche de minha trisavó se
torna --- naquele momento --- seu mundo por completo. E finalmente,
neste imenso oceano de apagamentos e refazeres, talvez --- agora ---
seja a hora de partir novamente do começo, e dessa vez com mais
distância.

\begin{center}\adforn{49}\end{center}

\noindent{}Existe uma história talmúdica que diz que o mundo foi criado e
destruído 26 vezes. O mundo já nasceu da e na ruína. Rabi
Bunam\footnote{Rabi Bunam de Pzhysha (1765--1827) foi um dos principais
  líderes do movimento \emph{hassídico} na Polônia durante o final do
  século \textsc{xviii} e o início do século \textsc{xix}. Seus contos, provérbios e
  histórias foram coletados e compilados no livro \emph{História do
  Rabi} do historiador e filósofo Martin Buber.} costumava dizer:
``No dia do Ano Novo\footnote{\emph{Rosh Hashaná} (em português,
  ``Cabeça do Ano'') é o dia festivo de início do calendário judaico.
  Acontece no primeiro dia do mês judaico de Tishrei, que corresponde
  normalmente a finais de setembro e inícios de outubro no calendário
  gregoriano.} o mundo começa novamente e, antes de começar,
termina. Assim como, antes de morrer, todas as forças agarram-se à vida,
o homem, no dia do Ano Novo, deve agarrar-se à vida do mundo com todas
as forças''\footnote{\textsc{buber}, \emph{op.\,cit.}, p.\,577.}.

A ``ruína'' e o ``agarrar-se à vida'' que nestas páginas
tenta-se construir --- por meio do esmiuçamento das territorialidades
temporárias, das línguas, da diáspora --- estão presentes nos documentos
e nos arquivos coletados e armazenados pelo grupo judaico \emph{Oyneg
Shabes}\footnote{\emph{Oyneg Shabes} foi um grupo judaico clandestino
  liderado pelo historiador polonês Emanuel Ringelblum (1900--1944) que
  atuou entre os anos de 1939 e 1943 dentro dos limites do gueto de
  Varsóvia, Polônia. O grupo dedicava-se a coletar e preservar relatos,
  histórias, sonhos, contos, gritos, choros e demais testemunhos sobre a
  vida individual e comunitária no gueto de Varsóvia durante a Segunda
  Guerra Mundial. Algumas horas antes do levante do gueto, em 1943, o
  grupo enterrou seus arquivos em um ato de esperança de que esses
  documentos fossem encontrados ao final da guerra e que suas memórias
  pudessem, ao menos, ser preservadas.} --- que dedicava-se a escavar
cada centímetro da história vivida entre 1939 e 1943 dentro do gueto de
Varsóvia. Seus registros são uma recolha incansável e imensurável de
vida (e de morte) e têm por intuito conservar migalhas de histórias,
mesmo daquelas mais particulares\footnote{A noção de particulares, que
  aqui tenta-se construir, está ligada ao entendimento de que as
  famílias, as comunidades e os indivíduos dentro do gueto possuíam suas
  próprias histórias e trajetórias de vida. Mesmo assim, suas histórias
  singulares, ao final, são formadas por meio de papéis amontoados que
  ressoam, por mais diferentes que sejam, em uma história coletiva.}.

Em \emph{Cosmopoéticas do refúgio}, o filósofo Dénètem Touam Bona, sobre
a identidade fronteiriça e a tentativa de construção de mundo que, ainda
hoje, torce a ``linha da cor'', diz:

\begin{quote}
Sobreviver num universo em ruína é se deixar atravessar, deixar-se
habitar pelo acréscimo da vida prodigado pelos ancestrais, pelos
animados (animais, plantas e povos do infinitamente pequeno) e pelos
elementos: abraçar a própria morte, a potência da sombra e do húmus,
para renascer\footnote{\textsc{bona}, \emph{Cosmopoéticas do refúgio}, p.\,85-86.}.
\end{quote}

Os testemunhos do grupo \emph{Oyneg Shabes}, essa espécie de
``{[}\ldots{}{]} abraçar a própria morte {[}\ldots{}{]} para
renascer''\footnote{\emph{Ibidem}.}, que por longo tempo permaneceram sob os
escombros do gueto de Varsóvia --- soterrados de neve e de terra ---,
manifestam-se não somente enquanto uma busca por uma história
particular, mas também enquanto uma narrativa montada com papéis
desatados por cada uma das histórias e das tragédias singulares. Este é
o caso de minha trisavó Dvora e de milhares de outras pessoas fugitivas
de guerra, que deixaram suas famílias em busca de promessas de
vidas\footnote{E pensar, o que elas diriam sobre Gaza?}. Afinal, como
nos lembra Didi-Huberman\footnote{\textsc{didi-huberman}, \emph{op.\,cit.}, p.\,107.}, o
paradoxo está em entender que judeus e judias foram unidos em morte, mas
que, em vida, viveram esparsamente: atravessados por brechas e fissuras.

Na falta de outras verdades anseio em acreditar que os parentes de minha
trisavó estavam unidos, pelo papel e pela carne, entre aqueles judeus e
judias que se levantaram contra o regime nazista durante o levante do
gueto de Varsóvia em 1943; e quiçá suas memórias, seus
``papéis-desejos'', seus ``papéis quaisquer'' e seus
``papéis chorados'' possam, ao menos, ter sido coletados,
armazenados e enterrados pelo grupo \emph{Oyneg Shabes}. E que assim
suas memórias, seus registros, suas histórias --- mesmo que eu não os
possua --- possam, de alguma maneira, perdurar no Tempo.

Os cunhados-irmãos\footnote{Após a morte de seu marido, Dvora foge da
  Polônia em direção aos Estados Unidos da América, deixando seu filho,
  meu bisavô Aba (Alberto) Najszteter, aos cuidados de seus cunhados e
  de seus sogros. A expressão ``cunhados-irmãos'' quer traduzir a
  proximidade que Dvora teria com os irmãos e irmãs de seu falecido
  marido.} de minha trisavó Dvora\footnote{Aqui é a nossa obrigação
  dar-lhes nome: Moshe Mekhl Najszteter \textsc{z''l}, Beila Myriam
  Sherberg \textsc{z''l}, Icchak Leib Najszteter \textsc{z''l}, Yehudit
  Sheindla Sukholski \textsc{z''l}, Mordechai Yossef Najszteter
  \textsc{z''l}, Hanna Sifra Zilbershpan \textsc{z''l}, Tova Shifra Godnizki
  \textsc{z''l}.} foram sistematicamente enviados ao campo de concentração
e de extermínio de Treblinka, na Polônia, a uma hora e meia de trem da
cidade de Varsóvia. Assim como foram os mais de 300 mil judeus e judias
residentes no gueto de Varsóvia --- dos quais seus cunhados-irmãos
também faziam parte --- e que se levantaram contra o regime nazista
entre abril e maio de 1943. Dos testemunhos-registros dos
cunhados-irmãos de minha trisavó Dvora me restam alguns nomes,
sobrenomes e longínquos graus de parentesco --- possivelmente alguma
carga genética --- o fio genealógico\footnote{A preservação genealógica
  --- seja ela sanguínea, textual ou cultural --- é uma das mais
  importantes preservações da vida das comunidades judaicas nas
  diásporas. Segundo o escritor Amós Oz e a historiadora Fania Oz-Salzberger:
  ``Para ser judeu não é necessário ser \emph{filho/\,filha de mãe judia},
  basta apenas ser um leitor''. Isso porque, na tradição judaica, todo
  leitor é revisor de originais, todo aluno é crítico e todo escritor,
  inclusive o criador do Universo, incorre em um grande número de
  questões.}. Dos milhares de judeus residentes do gueto de Varsóvia nos
restam resquícios de histórias, sonhos, canções,
gritos, ``papéis-desejos'', promessas de retorno, de persistência,
de vida. Dos judeus e judias perseguidos e mortos durante o Holocausto
nós mantivemos --- conservamos e fazemos perdurar --- ínfimos impulsos
de sobrevivência, gestos de resistência, fugas, resiliências e
sobrevidas, hinos partisanos e ``a lamentação para que nos
ensine, nos levante''\footnote{\textsc{didi-huberman}, \emph{op.\,cit.}, p.\,129.}.

Hoje, Treblinka --- este campo despovoado\footnote{Os nazistas
  destruíram não apenas os lugares judaicos de culto, de reza, de vida e
  de comunidade, como também seus lugares de morte.} e de morte, sítio
arqueológico, e possivelmente um suposto memorial --- se preenche de
deformadas pedras por toda a sua extensão. São formações rochosas
gravadas --- em nome\footnote{A morte pode possuir muitos nomes. Nas
  mais diversas pedras, espalhadas ao longo do memorial-campo de
  Treblinka, estão gravados nomes de lugares de memória que jamais
  devemos nos esquecer.} e em carne --- numa terra infértil. São elas,
as próprias pedras, memórias eternas, faíscas das chamas das almas de
pessoas mortas, que, espraiadas pelo ar, e levadas pelo velho vento do
sopro nômade da diáspora, se tornam fagulhas em que suas recordações e
memórias persistem e perduram no tempo, permanecem e criam
territorialidades temporárias por meio da língua, da conservação da
tradição, do indeterminismo da diáspora, e assim, são levadas de geração
em geração. São as faíscas das almas as próprias memórias que retornam;
e como diria Rabi Bunam, que antes da ruína agarram-se à vida.

Na tradição judaica há o costume de assentar pedras em túmulos nos
cemitérios como manifestação de recordo, de presença e da crença na
memória permanente. Certamente a materialidade de uma pedra não é
infinita, porém quiçá seja ela eterna em um determinado espaço-tempo ---
e, assim como a complexidade de Deus, essa energia não é infinita porque
não tem fim, antes é eterna porque não tem começo\footnote{Como diria
  Gabriel García Márquez: ``É a vida, mais do que a morte, que não
  tem limites''.}.

Após a guerra, foram encontradas, ao todo, 35.369 páginas esparsas
enterradas nos subsolos do gueto de Varsóvia pelo \emph{Oyneg Shabes}.
Este total representa apenas dois terços do material enterrado pelo
grupo. O restante --- um terço do material --- jaz --- afogado em terra
e em lama --- nas profundezas fúnebres de um solo gélido. Curiosamente,
o ``escovar a História a contrapelo''\footnote{``Escovar a
  história a contrapelo'' é uma metáfora de Walter Benjamin que busca
  analisar o passado questionando e revisitando a narrativa oficial dos
  vencedores. Por meio dessa abordagem, foca-se nas experiências dos
  oprimidos, vencidos e silenciados, revelando as lacunas e resistências
  que o progresso linear ignorou.}, sobre que o filósofo Walter Benjamin
há tempos já discorria, é literalmente materializado nos escombros do
gueto. São pedras, tijolos, pedaços de concreto, estilhaços de vidro e
demais resíduos da construção civil levantados e retirados em busca de
testemunhos, de folhas sujas, amassadas e soterradas --- e que, por fim,
sobrevivem às intempéries da História e do Tempo.

% Comentário da editora: Acho que está tudo bem usar história e tempo em caixa baixa aqui. Existem também outras ocasiões em que termos parecidos estão em caixa alta sem necessidade. Posso padronizar?

Usando novamente da deposição do filósofo Dénètem Touam Bona: ``As
guerrilhas se apresentam como uma não batalha, um disfarce e uma
camuflagem que zombam com a moral dos poderosos''\footnote{\textsc{bona}, op.
  cit., p.\,41.}. Dessa maneira, o grupo \emph{Oyneg Shabes}, para além
de uma tentativa de preservação histórica, brilhantemente entende que,
para o ``povo do livro'', basta também que suas páginas
sobrevivam, que seus traços de humanidade possam ser lidos em letras
sobre papéis. É por meio do encontro entre o fino risco da caneta e a
matéria que a palavra transforma-se em lugar. É a partir destes gestos
que recordamos os 6 milhões de judeus e judias mortos durante o
Holocausto --- e assim, compreendemos que suas almas retornam e
retornarão, agarrando-se à vida --- perdurando em cada um dos finos
traços de letras.

\begin{center}\adforn{49}\end{center}

Em março de 1943, no gueto de Vilnius\footnote{O gueto de Vilnius,
  na Lituânia, durou entre os anos 1941--1943.}, o poeta Avrom Sutzkever escreveu um poema. Através de suas palavras e imagens de sobrevivência, são os \emph{papéis-testemunhos}, mais do que a própria violência, que perduram. A materialidade frágil de uma folha de papel pode ser amassada, rasgada, transformada em cinzas. Mas são nas palavras, escritas ou faladas, que as histórias permanecem vivas e resgatam sentidos de resistência, atualizados a cada geração. E o momento de destruição dá ``{[}\ldots{}{]} à luz a uma estrela dançarina''\footnote{\textsc{bey}, \emph{\textsc{taz} --- Zona Autônoma Temporária}, p.,7.}. 

% Comentário da editora: Cortei essa parte.
% um esforço para o surgimento de desejos latentes de semear, uma natureza de migalhas a serem coletadas e, por fim, espalhadas como finos lapsos de fagulhas e centelhas.

\begin{verse}
Talvez essas palavras também\\
Vão durar, e chegada a hora,\\
Subir para a luz\\
E florescer inesperadamente\\
E tal como o antigo grão\\
Tornou-se espiga\\
Talvez estas palavras vão nutrir,\\
Talvez essas palavras vão permanecer\\
Ao povo, em seu incessante caminho.\footnote{\textsc{sutzkever} \emph{apud}
  \textsc{didi-huberman}, \emph{op.\,cit.}, p.\,128.}
\end{verse} 

% Comentário da editora: Cortei pra não ficar pulando de um assunto pro outro.
% Desde as cinzas das chamas das fogueiras inquisitórias aos papéis dos guetos, 

São estes os registros das odisseias particulares e coletivas, de corpos individuais e comunitários que percorrem e ecoam os movimentos elipsoidais das diásporas, nascendo e se construindo em territórios temporários, passíveis de apropriação e pertencimento. É hora de observar, de perto e de longe, o que brotará desses outros locais.\footnote{Pensar um judaísmo por uma ótica que considere as particularidades de uma identidade, de uma cultura e de uma localidade que se encontra (neste momento) em um novo processo de adaptação, resiliência, câmbio e debate interno.} \emph{Mir zaynen do}.\footnote{Do iídiche, ``Nós estamos aqui''. É também a expressão entoada no refrão do \emph{Hino dos Partisanos}.}

% \begin{center}\adforn{49}\end{center}

% \noindent{}Ainda prefiro a versão das coisas em que não se sabe ao certo o que foi
% lido, o que foi inventado. Se diz que existem apenas sete possíveis
% histórias a serem contadas no mundo, e o que muda entre cada uma delas
% seria a capacidade de rearranjo de cada coletor, de cada recorte, de
% cada ritmo, de cada jeito de entrelaçar essas narrativas. As sete
% histórias não saberia dizer, porém sei que continuamos a escrever sempre
% o mesmo livro.

% Em \emph{A ensaificação de tudo}, a escritora norte-americana Christy
% Wampole diz que ``O ensaísta está interessado em pensar sobre si
% pensando sobre as coisas''\footnote{\textsc{wampole}, \emph{A ensaificação de
%   tudo}.}. Ao fim e ao cabo, esta pesquisa é uma espécie de ensaio que,
% assim como a própria diáspora, sustenta-se enquanto tentativa de
% solução, ainda que na ausência de soluções. São impulsos e propostas
% para ``pensar sobre si pensando sobre as coisas''\footnote{\emph{Ibidem}.},
% e dessa maneira, construir outras territorialidades e frequências
% comunitárias. Assim sendo:

% \begin{quote}
% Esse não é um livro de história. A escolha que nele se encontrará não
% seguiu outra regra mais importante do que meu gosto, meu prazer, uma
% emoção, o riso, a surpresa, um certo assombro ou qualquer outro
% sentimento, do qual teria dificuldades, talvez, em justificar a
% intensidade, agora que o primeiro momento da descoberta
% passou\footnote{\textsc{foucault}, \emph{A vida dos homens infames}, p.\,203.}.
% \end{quote}

% Busca-se aqui entender que as mais diversas \emph{perspectivas
% judaico-diaspóricas sobre o lugar} estão entremeadas e entrelaçadas de
% signos, símbolos, línguas e arquiteturas --- em suas mais variadas
% escalas --- em que se permite o nascimento e a consolidação de
% identidades fluidas, regadas de indeterminabilidade e, por isso,
% passíveis de apropriação e de pertencimento.

% São as fagulhas, as chispas, as faíscas, as memórias-palavras que
% perduram, que vagueiam pelo ar e são levadas pelo velho sopro do vento
% nômade da diáspora. E, ao pousarem, criam territorialidades que, mesmo
% que temporárias, possibilitam a construção de localidades próprias e de
% sentidos de identificação. Dessa maneira, os lugares que aqui procura-se
% trabalhar, mais do que espacialidades físicas, se manifestam em seus
% sentidos identitários e comunitários e através das línguas, dos objetos
% culturais e de ritos religiosos. São eles: os espaços de trânsito, do
% indeterminado, do fluido e do híbrido que cruzam --- em suas mais
% diversas maneiras --- as noções espaço-temporais judaicas.

% Essas geografias de identidades apropriadas serão esmiuçadas e
% investigadas no decorrer deste texto por meio dos conceitos judaicos de:

% \begin{enumerate}
% \item \textsc{bs'd}; 
% \item Kipá; 
% \item Mezuzá; 
% \item Eruv; 
% \item Makom. 
% \end{enumerate}

% Transversalmente, procura-se entender
% a página, as palavras, as línguas diaspóricas\footnote{As línguas também
%   possuem biografia. O ladino conserva as memórias de infância de
%   um espanhol medieval. As pessoas podem esquecer, mas as línguas não.
%   As línguas mantêm as alegrias, as dores e as feridas de tudo que há
%   passado.} como elementos essenciais para a criação de corpos
% identitários e comunitários; a \emph{Shechiná} --- o lugar de habitação,
% de crença e de presença divina --- como uma territorialidade temporária;
% e a diáspora enquanto um limite e um limiar em si, um espaço eternamente
% transitório.

% Rabi Tarfon\footnote{Rabi Tarfon foi membro da terceira geração de
%   estudiosos da \emph{Mishná}, vivendo entre os séculos \textsc{i} e \textsc{ii} d.e.c.}
% costumava dizer: ``Não lhe é exigido que complete a tarefa, mas
% você não é livre para escapar dela''\footnote{\emph{Pirkei Avot} 2:16,
%   escrito cerca do século \textsc{i}.}. Essa constante elaboração é o que
% impossibilita o esgotamento de sentidos sobre o ser judeu. Por fim, já
% que estaremos sempre a escrever o mesmo livro, essa investigação
% sustenta-se enquanto um ensaio identitário, sem a pretensão de alicerçar
% e fixar imagens judaicas.

% Em resumo, esta breve, dispersa e pequena introdução é uma abertura de
% portas.

% \emph{I ke io lo mire}\footnote{Do ladino: ``e que eu o veja''.}.